\documentclass[english,twoside, openright, hidelinks, 11 pt, class=report,crop=false]{standalone}
\usepackage[utf8]{inputenc}
\usepackage{lmodern} % load a font with all the characters
\usepackage{geometry}
\geometry{verbose,paperwidth=16 cm, paperheight=24 cm, inner=2.3cm, outer=1.7 cm, bmargin=2cm, tmargin=1.8cm}

\setlength{\parindent}{0bp}
\usepackage{import}
\usepackage[subpreambles=false]{standalone}
\usepackage{amsmath}
\usepackage{amssymb}
\usepackage{esint}
\usepackage{babel}
\usepackage{tabu}
\makeatother
\makeatletter

\usepackage{titlesec}
\usepackage{ragged2e}
\RaggedRight
\raggedbottom
\frenchspacing

\usepackage{graphicx}
\usepackage{float}
\usepackage{subfig}
\usepackage{placeins}
\usepackage{cancel}
\usepackage{framed}
\usepackage{wrapfig}
\usepackage[subfigure]{tocloft}
\usepackage[font=footnotesize,labelfont=sl]{caption} % Figure caption
\usepackage{bm}
\usepackage[dvipsnames, table]{xcolor}
\definecolor{shadecolor}{rgb}{0.105469, 0.613281, 1}
\definecolor{blue}{cmyk}{1, 0, 0, 0}
\definecolor{red}{cmyk}{0, 1, 1, 0}
\definecolor{green}{cmyk}{100,0,100,0}
\colorlet{shadecolor}{Emerald!15} 
\usepackage{icomma}
\makeatother
\usepackage[many]{tcolorbox}
\usepackage{multicol}
\usepackage{stackengine}
\usepackage{listings}
\usepackage{esvect} %For vectors with capital letters

% For tabular
\usepackage{array}
\usepackage{multirow}
\usepackage{longtable} %breakable table

% Ligningsreferanser
\usepackage{mathtools} % for mathclap

%\mathtoolsset{showonlyrefs}

% sections without numbering in toc
\newcommand\tsec[1]{\phantomsection \addcontentsline{toc}{section}{#1}
	\section*{#1}}

% index
\usepackage{imakeidx}
\makeindex[title=Indeks]

%Footnote:
\usepackage[bottom, hang, flushmargin]{footmisc}
\usepackage{perpage} 
\MakePerPage{footnote}
\addtolength{\footnotesep}{2mm}
\renewcommand{\thefootnote}{\arabic{footnote}}
\renewcommand\footnoterule{\rule{\linewidth}{0.4pt}}
\renewcommand{\thempfootnote}{\arabic{mpfootnote}}

%colors
\definecolor{c1}{cmyk}{0,0.5,1,0}
\definecolor{c2}{cmyk}{1,0.25,1,0}
\definecolor{n3}{cmyk}{1,0.,1,0}
\definecolor{neg}{cmyk}{1,0.,0.,0}


\newcommand{\nreq}[1]{
\begin{equation}
	#1
\end{equation}
}


% Equation comments
\newcommand{\cm}[1]{\llap{\coln{ #1}}}


\usepackage[inline]{enumitem}
\newcounter{rg}
\numberwithin{rg}{chapter}

\newcommand{\reg}[2][]{\begin{tcolorbox}[boxrule=0.3 mm,arc=0mm,colback=Blue!3] {\refstepcounter{rg}\phantomsection \large \textbf{\therg \;#1} \vspace{5 pt}}\newline #2  \end{tcolorbox} \vs}
\newcommand{\regdef}[2][]{\begin{tcolorbox}[boxrule=0.3 mm,arc=0mm,colback=Blue!3] {\refstepcounter{rg}\phantomsection \large \textbf{\therg \;#1} \vspace{5 pt}}\newline #2  \end{tcolorbox} \vs}
\newcommand{\words}[1]{\begin{tcolorbox}[boxrule=0.3 mm,arc=0mm,colback=Teal!3] #1  \end{tcolorbox} \vs}

\newcommand\alg[1]{\begin{align*} #1 \end{align*}}

\newcommand\eks[2][]{\begin{tcolorbox}[boxrule=0.3 mm,arc=0mm,enhanced jigsaw,breakable,colback=Green!5] {\large \textbf{\ekstitle #1} \vspace{5 pt}\\} #2 \end{tcolorbox}\vs}

\newcommand{\st}[1]{\begin{tcolorbox}[boxrule=0.0 mm,arc=0mm,enhanced jigsaw,breakable,colback=Yellow!7]{ #1} \end{tcolorbox}}

\newcommand{\spr}[1]{\begin{tcolorbox}[boxrule=0.3 mm,arc=0mm,enhanced jigsaw,breakable,colback=Yellow!7] {\large \textbf{\sprtitle} \vspace{5 pt}\\} #1 \end{tcolorbox}\vs}

\newcommand{\sym}[1]{\colorbox{NavyBlue!15}{#1}}

\newcommand{\info}[2]{\begin{tcolorbox}[boxrule=0.3 mm,arc=0mm,enhanced jigsaw,breakable,colback=Cyan!6] {\large \textbf{#1} \vspace{5 pt}\\} #2 \end{tcolorbox}
}

\newcommand\algv[1]{\vspace{-11 pt}\begin{align*} #1 \end{align*}}

\newcommand{\regv}{\vspace{5pt}}
\newcommand{\mer}{\textsl{\note}: }
\newcommand{\mers}[1]{{\footnotesize \mer #1}}
\newcommand\vsk{\vspace{11pt}}
\newcommand{\tbs}{\vspace{5pt}}
\newcommand\vs{\vspace{-11pt}}
\newcommand\vsb{\vspace{-16pt}}
\newcommand\vso{\vspace{-5pt}}
\newcommand\br{\\[5 pt]}
\newcommand{\figp}[1]{../fig/#1}
\newcommand\algvv[1]{\vs\vs\begin{align*} #1 \end{align*}}
\newcommand{\y}[1]{$ {#1} $}
\newcommand{\os}{\\[5 pt]}
\newcommand{\prbxl}[2]{
\parbox[l][][l]{#1\linewidth}{#2
	}}
\newcommand{\prbxr}[2]{\parbox[r][][l]{#1\linewidth}{
		\setlength{\abovedisplayskip}{5pt}
		\setlength{\belowdisplayskip}{5pt}	
		\setlength{\abovedisplayshortskip}{0pt}
		\setlength{\belowdisplayshortskip}{0pt} 
		\begin{shaded}
			\footnotesize	#2 \end{shaded}}}
\newcommand{\fgbxr}[2]{
	\parbox[r][][l]{#1\linewidth}{#2
}}		

\renewcommand{\cfttoctitlefont}{\Large\bfseries}
\setlength{\cftaftertoctitleskip}{0 pt}
\setlength{\cftbeforetoctitleskip}{0 pt}

\newcommand{\bs}{\\[3pt]}
\newcommand{\vn}{\\[6pt]}
\newcommand{\fig}[1]{\begin{figure}[H]
		\centering
		\includegraphics[]{\figp{#1}}
\end{figure}}

\newcommand{\localfig}[2]{\begin{figure}[H]
		\centering
		\includegraphics[scale=#2]{#1}
\end{figure}}

\newcommand{\figc}[2]{\begin{figure}
		\centering
		\includegraphics[]{\figp{#1}}
		\caption{#2}
\end{figure}}
\newcommand{\arc}[1]{{
		\setbox9=\hbox{#1}%
		\ooalign{\resizebox{\wd9}{\height}{\texttoptiebar{\phantom{A}}}\cr\textit{#1}}}}

\newcommand{\sectionbreak}{\clearpage} % New page on each section

\newcommand{\nn}[1]{
\begin{equation*}
	#1
\end{equation*}
}

\newcommand{\enh}[1]{\,\text{#1}}

%asin, atan, acos
\DeclareMathOperator{\atan}{atan}
\DeclareMathOperator{\acos}{acos}
\DeclareMathOperator{\asin}{asin}

% Comments % old cm, ggb cm is new
%\newcommand{\cm}[1]{\llap{\color{blue} #1}}

%%%

\newcommand\fork[2]{\begin{tcolorbox}[boxrule=0.3 mm,arc=0mm,enhanced jigsaw,breakable,colback=yellow!7] {\large \textbf{#1 (\expl)} \vspace{5 pt}\\} #2 \end{tcolorbox}\vspace{-5pt} }
 
%colors
\definecolor{projectgreen}{cmyk}{100,0,100,0}

\newcommand{\colr}[1]{{\color{Red} #1}}
\newcommand{\colb}[1]{{\color{Cyan} #1}}
\newcommand{\colo}[1]{{\color{Orange} #1}}
\newcommand{\colc}[1]{{\color{SeaGreen} #1}}
\newcommand{\colg}[1]{{\color{Green} #1}}
\newcommand{\coln}[1]{{\color{NavyBlue} #1}}

% Methods
\newcommand{\metode}[2]{
	\textsl{#1} \\ \vspace{-10pt}
	\rule{#2}{0.75pt}
}

%Opg
\newcommand{\abc}[1]{
	\begin{enumerate}[label=\alph*),leftmargin=18pt]
		#1
	\end{enumerate}
}
\newcommand{\abcs}[2]{
	\begin{enumerate}[label=\alph*),start=#1,leftmargin=18pt]
		#2
	\end{enumerate}
}
\newcommand{\abcn}[1]{
	\begin{enumerate}[label=\arabic*),leftmargin=18pt]
		#1
	\end{enumerate}
}




\newcommand{\abch}[1]{
	\hspace{-2pt}	\begin{enumerate*}[label=\alph*), itemjoin=\hspace{1cm}]
		#1
	\end{enumerate*}
}

\newcommand{\abchs}[2]{
	\hspace{-2pt}	\begin{enumerate*}[label=\alph*), itemjoin=\hspace{1cm}, start=#1]
		#2
	\end{enumerate*}
}

% Exercises


\newcounter{opg}
\numberwithin{opg}{section}

\newcounter{grub}
\numberwithin{opg}{section}
\newcommand{\op}[1]{\vspace{15pt} \refstepcounter{opg}\large \textbf{\color{NavyBlue}\theopg} \vspace{2 pt} \label{#1} \\}
\newcommand{\opd}[1]{\vspace{15pt} \refstepcounter{opg}\large \textbf{\color{NavyBlue}\theopg} \includegraphics{../../calculatorimage} \vspace{2 pt} \label{#1} \\}
\newcommand{\eksop}[2]{\vspace{15pt} \refstepcounter{opg}\large \textbf{\color{NavyBlue}\theopg} (#1) \vspace{2 pt} \label{#2} \\}

\newcommand{\nes}{\stepcounter{section}
	\setcounter{opg}{0}}
\newcommand{\opr}[1]{\vspace{3pt}\textbf{\ref{#1}}}
\newcommand{\oeks}[1]{\begin{tcolorbox}[boxrule=0.3 mm,arc=0mm,colback=white]
		\textit{\ekstitle: } #1	  
\end{tcolorbox}}
\newcommand\opgeks[2][]{\begin{tcolorbox}[boxrule=0.1 mm,arc=0mm,enhanced jigsaw,breakable,colback=white] {\footnotesize \textbf{\ekstitle #1} \\} \footnotesize #2 \end{tcolorbox}\vspace{-5pt} }


% tag exercises
\newcommand{\tagop}[1]{ 
{\small \color{Gray} #1} \os
}

% License
\newcommand{\lic}{
This book is part of the \net{https://sindrsh.github.io/openmathbooks/}{OpenMathBooks} project. OpenMathBooks © 2022 by Sindre Sogge Heggen is licensed under CC BY-NC-SA 4.0. To view a copy of this license, visit \net{http://creativecommons.org/licenses/by-nc-sa/4.0/}{http://creativecommons.org/licenses/by-nc-sa/4.0/}}

%referances
\newcommand{\net}[2]{{\color{NavyBlue}\href{#1}{#2}}}
\newcommand{\hrs}[2]{\hyperref[#1]{\color{NavyBlue}#2 \ref*{#1}}}
\newcommand{\refunnbr}[2]{\hyperref[#1]{\color{NavyBlue}#2}}


\newcommand{\openmath}{\net{https://sindrsh.github.io/openmathbooks/}{OpenMathBooks}}
\newcommand{\am}{\net{https://openmath.online/books/pdf/AM_bm.pdf}{AM1}}
\newcommand{\mb}{\net{https://openmath.online/books/pdf/MB_bm.pdf}{MB}}
\newcommand{\tmen}{\net{https://openmath.online/books/pdf/TM1_bm.pdf}{TM1}}
\newcommand{\tmto}{\net{https://openmath.online/books/pdf/TM2_bm.pdf}{TM2}}
\newcommand{\amto}{\net{https://sindrsh.github.io/FirstPrinciplesOfMath/}{AM2}}
\newcommand{\eksbm}{
\footnotesize
Dette er opppgaver som har blitt gitt ved sentralt utformet eksamen i Norge. Oppgavene er laget av Utdanningsdirektoratet. Forkortelser i parantes viser til følgende:
\begin{center}
	\begin{tabular}{c|c}
		E & Eksempeloppgave \\
		V/H & Eksamen fra vårsemesteret/høstsemesteret\\
		G/1P/1T/R1/R2 & Fag  \\
		XX & År 20XX \\
		D1/D2 & Del 1/Del 2
	\end{tabular}
\end{center}
Tekst og innhold kan her være noe endret i forhold til originalen.
}

%Excel og GGB:

\newcommand{\g}[1]{\begin{center} {\tt #1} \end{center}}
\newcommand{\gv}[1]{\begin{center} \vspace{-11 pt} {\tt #1}  \end{center}}
\newcommand{\cmds}[2]{{\tt #1}\\
	#2}

% outline word
\definecolor{Teal}{cmyk}{1, 0, 0.4, 0.3}
\newcommand{\outl}[1]{{\boldmath \color{Teal}\textbf{#1}}}
%line to seperate examples
\newcommand{\linje}{\rule{\linewidth}{1pt} }


%Vedlegg
\newcounter{vedl}
\newcounter{vedleq}
\renewcommand\thevedl{\Alph{vedl}}	
\newcommand{\nreqvd}{\refstepcounter{vedleq}\tag{\thevedl \thevedleq}}

%%% Writing code

\definecolor{codegreen}{rgb}{0,0.6,0}
\definecolor{codegray}{rgb}{0.5,0.5,0.5}
\definecolor{codepurple}{rgb}{0.58,0,0.82}
\definecolor{backcolour}{rgb}{0.95,0.95,0.92}


\lstdefinestyle{mystyle}{
	backgroundcolor=\color{backcolour},   
	commentstyle=\color{codegreen},
	keywordstyle=\color{magenta},
	numberstyle=\tiny\color{codegray},
	stringstyle=\color{codepurple},
	basicstyle=\ttfamily\footnotesize,
	breakatwhitespace=false,         
	breaklines=true,                 
	captionpos=b,                    
	keepspaces=true,                 
	numbers=left,                    
	numbersep=5pt,                  
	showspaces=false,                
	showstringspaces=false,
	showtabs=false,                  
	tabsize=2
}

\lstset{literate=%
	{æ}{{\ae}}1
	{ø}{{\o}}1
	{å}{{\aa}}1
}

\lstset{style=mystyle}


\newcommand{\pymet}[1]{{\ttfamily\color{magenta} #1}}
\newcommand{\pytype}[1]{{\ttfamily\color{codepurple} #1}}

\newcommand{\python}[1]{
\begin{tcolorbox}[boxrule=0.3 mm,arc=0mm,colback=white]
\lstinputlisting[language=Python]{#1}
\end{tcolorbox}}
\newcommand{\pythonut}[2]{
\begin{tcolorbox}[boxrule=0.3 mm,arc=0mm,colback=white]
\small 
%\textbf{Kode}
\lstinputlisting[language=Python]{#1}	
\vspace{11pt}
\textbf{Utdata} \\ \ttfamily
#2
\end{tcolorbox}}
%%%

%page number
%\usepackage{fancyhdr}
%\pagestyle{fancy}
%\fancyhf{}
%\renewcommand{\headrule}{}
%\fancyhead[RO, LE]{\thepage}

\newcommand{\deltax}{\,\text{d}x}
\newcommand{\deltau}{\,\text{d}u}
\usepackage{datetime2}
%%\usepackage{sansmathfonts} for dyslexia-friendly math
\usepackage[]{hyperref}

\newcommand{\texandasy}{Teksten er skrevet i \LaTeX\ og figurene er lagd ved hjelp av Asymptote.}

\newcommand{\expl}{forklaring}

\newcommand{\note}{Merk}
\newcommand{\notesm}[1]{{\footnotesize \textsl{\note:} #1}}
\newcommand{\ekstitle}{Eksempel }
\newcommand{\sprtitle}{Språkboksen}
\newcommand{\vedlegg}[1]{\section*{Vedlegg \thevedl: #1}  \setcounter{vedleq}{0} \refstepcounter{vedl}}

\newcommand\sv{\vsk \textbf{Svar} \vspace{4 pt}\\}

% exercises
\newcommand{\opgt}{\newpage \phantomsection \addcontentsline{toc}{section}{Oppgåver} \section*{Oppgaver for kapittel \thechapter}\vs \setcounter{section}{1}}

\newcommand{\grubop}[1]{\vspace{15pt} \refstepcounter{grub}\large \textbf{\color{NavyBlue} Gruble \thegrub} \vspace{2 pt} \label{#1} \\}
\newcommand{\grubeksop}[2]{\vspace{15pt} \refstepcounter{grub}\large \textbf{\color{NavyBlue} Gruble \thegrub} (#2) \vspace{2 pt} \label{#1} \\}
\newcommand{\gruboprecc}[1]{\vspace{15pt} \refstepcounter{grub}\large \textbf{\color{NavyBlue} Gruble* \thegrub} \vspace{2 pt} \label{#1} \\}
\newcommand{\grubeksoprecc}[2]{\vspace{15pt} \refstepcounter{grub}\large \textbf{\color{NavyBlue} Gruble* \thegrub} (#2) \vspace{2 pt} \label{#1} \\}
\newcommand{\grubr}[1]{\vspace{3pt}\textbf{Gruble \ref{#1}}}

%references
\newcommand{\reftab}[1]{\hrs{#1}{tabell}}
\newcommand{\rref}[1]{\hrs{#1}{regel}}
\newcommand{\dref}[1]{\hrs{#1}{definisjon}}
\newcommand{\refkap}[1]{\hrs{#1}{kapittel}}
\newcommand{\refsec}[1]{\hrs{#1}{seksjon}}
\newcommand{\refdsec}[1]{\hrs{#1}{delseksjon}}
\newcommand{\refved}[1]{\hrs{#1}{vedlegg}}
\newcommand{\eksref}[1]{\textsl{#1}}
\newcommand\fref[2][]{\hyperref[#2]{\textsl{figur \ref*{#2}#1}}}
\newcommand{\refop}[1]{{\color{blue}Oppgave \ref{#1}}}
\newcommand{\refops}[1]{{\color{blue}oppgave \ref{#1}}}

% Solutions manual
\newcommand{\selos}{Se løsningsforslag.}

\newcommand{\ompref}{Omgjøring av prefikser}


\begin{document}

\opgt

\op{opggjeromtilm}
Gjør om til antall meter. \os
\abch{
	\item 484\enh{km}
	\item 91\enh{km}
	\item 2\,402\enh{km}
}

\op{opggjeromtilg}
Gjør om til antall gram. \os
\abch{
	\item 484\enh{kg}
	\item 91\enh{hg}
	\item 2\,402\enh{hg}
}

\op{opggjeromtilL}
Gjør om til antall liter \os
\abch{
	\item 480\enh{dl}
	\item 9100\enh{cl}
	\item 24\,000\enh{cl}
}

\op{opggjeromblanda}
Gjør om 
\begin{multicols}{3}
	\abc{
		\item 12,4\enh{m} \\til antall km.
		\item 42\enh{dm} \\til antall m.
		\item 58,15\enh{cm} \\til antall mm.
		\item 0,0074\enh{km}\\ til antall m.
		\item 0,15\enh{m} \\til antall cm.
		\item 9,7\enh{g} \\til antall hg.
		\item 0,15\enh{mg}\\ til antall g.
		\item 1,419\enh{hg}\\ til antall mg.
		\item 31\enh{mg} \\til antall hg.
		\item 64\,039\enh{mg}\\ til antall kg.
		\item 89\enh{dl} \\til antall l.
		\item 691,4\enh{l}\\ til antall cl.
		\item 15\enh{l} \\til antall ml.
		\item 918\enh{cl} \\til antall l.
		\item 0,55\enh{dl} \\til antall ml.
	}
\end{multicols}

\nes


\op{opgvolprisl}
En prisme har lengde 9\enh{cm}, bredde 10\enh{cm} og høgde 8\enh{cm}.\os

Finn volumet til prismen.


\op{opgvolkjegl}
En kjegle har radius 6\enh{dm} og høgde 4\enh{dm}.
\abc{
	\item Finn volumet til kjeglen.
	\item Hvor mange liter rommer kjeglen?
}

\op{opgvolpyrl}
En firkantet pyramide har lengde 4\enh{cm}, bredde 9\enh{cm} og høgde 10\enh{cm}.
\abc{
	\item Finn volumet til kjeglen.
	\item Hvor mange liter rommer kjeglen?
}

\nes

\op{opgspeed}
\abc{
\item En bil kjører i 60\enh{km/h} i 3\enh{h}. Finn lengden bilen har kjørt.
\item En bil kjører i 50\enh{km/h} i 4,25\enh{h}. Finn lengden bilen har kjørt.
\item En bil kjører 375\enh{km} i løpet av 5\enh{h}. Hvilken fart har bilen?
}

\op{opgkilopris} \vs
\abc{
\item 10\enh{kg} tomater koster 35\enh{kr}. Hva er kiloprisen til tomatene?
}


\eksop{1PV26d1opg6}{1PV26D1}
Noen venner skal leie en seilbåt. Prisen hver person må betale er omvendt proporsjonal med antall personer som blir med. Fyll inn tallet som mangler i tabellen.
\begin{center}
	\begin{tabular}{|l|c|c|}
		\hline
		Antall personer & 5 & 14 \\ \hline
		Pris per person (i kroner) & 700 & \\ \hline
	\end{tabular}
\end{center}

\eksop{1ph22d2opg3}{1PH22D2}
For å gjøre om fra grader $F$ fahrenheit til grader $C$ celsius kan vi bruke
formelen
\[ C=\frac{F-32}{1,8} \]
Er grader celsius og grader fahrenheit proporsjonale størrelser, omvendt
proporsjonale størrelser, eller ingen av delene?

\op{opgstorlbolt}
Usain Bolt har verdensrekorden for 100\enh{m} sprint. I tabellen under ser du hva tidtakeren viste ved hver 10. meter under dette rekordløpet. \vs
\begin{center}\small
	\begin{tabular}{l|c|c|c|c|c|c|c|c|c|c}
meter & 10 & 20 & 30 & 40 & 50 & 60 & 70 & 80 & 90 & 100\\
sekunder& 1,89 & 2,88 & 3,78 & 4,64 & 5,47 & 6,29 & 7,10 & 7,92 & 8,75 & 9,58
\end{tabular}
\end{center}
\abc{
\item Hvis Bolts fart hadde vært den samme under hele løpet, hva hadde farten hans vært da?
\item Hvorfor er det ikke rimelig å anta at Bolt hadde den samme farten under hele løpet?
\item Anta at Bolt under dette løpet nådde den høgste farten et menneske har sprunget. Finn en tilnærming til denne farten.
}
\newpage
\nes
\op{opgstorlbolt2}
Gjør om svarene fra \refops{opgstorlbolt} a) og c) til hastigheter oppgitt i '\enh{km/h}'. 

\op{opgstorlfinn} 
Skriv ned eksempel på et dyr, et insekt, en gjenstand eller annet som veier mellom 1-100\,mg, cg, dg, g, dag, hg og kg.

\eksop{1PH23D1}{1PH23D1opg1}
Tobias lurer på hvor mye vann han bør drikke hver dag.\os
Han finner ulike svar på ulike nettsider.\\
På en nettside finner han teksten nedenfor.
\st{
Voksne har hvert døgn behov for ca. 30 mL væske per kilogram kroppsvekt.
Husk at vann er den beste tørstedrikken.
}
Tobias veier 70 kg. \os
Hvor mange liter vann bør Tobias drikke i løpet av et døgn ifølge nettsiden?

\grubop{1ph23d1opg3} (1PH23D1) \os
Ohms lov sier at strømmen $ I $ gjennom en metallisk leder med konstant temperatur
er proporsjonal med spenningen $ \Omega $ og omvendt proporsjonal med motstanden $ R $ i lederen.\os
Argumenter for om hver av påstandene er sann eller usann.\os
1) Hvis vi øker spenningen, vil strømmen også øke.\os
2) Hvis vi øker motstanden, vil strømmen også øke.

\grubeksoprecc{1pv21d1opg1}{1PV32D1}
En kopp med 220\,mL cappucino koster 22 kroner.	Hvor mye koster cappucionoen per liter?

\grubeksoprecc{1pv21d1opg8}{1PV21D1}
Vi regner at avstanden fra jorda til Saturn er $1,4\cdot10^9\enh{km}$. Et av romskipene som astronautene reiste med fra jorda til månen holdt en hastighet på $1,0\cdot10^6\enh{km/døgn}$. Hvor mange døgn ville det ta å reise fra jorda til Saturn med denne hastigheten?

\grubeksoprecc{1ph21d1opg4}{1PH21}
En tunell er 24 kilometer lang. En bil kjører gjennom halve tunnelen med en
gjennomsnittsfart på 80\,km/h. Resten av tunnelen kjører den med en gjennomsnittsfart på 60\,km/h. \os

Hvor mange minutter bruker bilen på å kjøre gjennom tunellen?

\gruboprecc{opgstorlpropeks}
\abc{
\item Gi et eksempel på to størrelser som er proporsjonale.
\item Gi et eksempel på to størrelser som er omvendt proporsjonale.
\item Skisser en graf som beskriver to proporsjonale størrelser $x$ og $y$. Hvor vil grafen skjære $y$-aksen?
\item Skisser en graf som beskriver to omvendt proporsjonale størrelser $x$ og $y$. Hvilken verdi vil $y$ nærme seg når $x$ øker i verdi?
}
Si at proporsjonalitetskonstanten er 20, og gitt tabellen
\begin{center}
	\begin{tabular}{l | c |c|}
		$x$ & 1 & 10 \\ \hline
		$y$ &   &  \\ 
	\end{tabular}
\end{center}
\abcs{5}{
\item Fyll inn $y$-verdiene i tabellen gitt at $x$ og $y$ er proporsjonale.
\item Fyll inn $y$-verdiene i tabellen gitt at $x$ og $y$ er omvendt proporsjonale.
}

\grubeksoprecc{1ph21d2opg4}{1PH21}
Avgjør hvilke av påstandene som er riktige.
\begin{center}
	\begin{tabular}{p{2cm} | p{8cm}}
		\hline
		Påstand 1 & Dersom utgiftene til en klassefest skal deles likt mellom elevene som er
		med på festen, vil beløpet hver elev må betale alltid være omvendt
		proporsjonalt med antall elever. \\ \hline
		Påstand 2 & To størrelser er alltid proporsjonale dersom det er slik at når den ene
		øker, så øker den andre også. \\ \hline
		Påstand 3& To størrelser er alltid omvendt proporsjonale dersom den ene størrelsen
		dobler seg når den andre halveres. \\ \hline
		Påstand 4& Arealet av en sirkel er alltid proporsjonalt med omkretsen av sirkelen.\\
		\hline	
	\end{tabular}
\end{center}

\grubeksoprecc{1pv23d2opg4}{1PV23D2}
En bonde selger sekker med poteter. I koordinatsystemet nedenfor ser du sammenhengen mellom vekt og pris for
potetsekkene. Hvert av punktene $A$, $B$, $C$, $D$, $E$ og $F$ representerer en potetsekk.
\fig{1pv23d2opg4}
\abc{
\item Hvilken sekk veier mest?
\item To av sekkene har lik pris. Hvilke to er det?
\item Hvilken av sekk $B$ og sekk $C$ har høyest kilopris?
\item Hvilke to sekker har lik kilopris?
}

\grubeksoprecc{1pv24d1opg2}{1PV24D1}
Under ser du grafene til fire funksjoner $f(x)$, $g(x)$, $p(x)$ og $q(x)$.
\abc{
\item Avgjør om en eller flere av grafene viser sammenhengen\\ mellom to størrelser
som er proporsjonale.
\item Avgjør om en eller flere av grafene viser sammenhengen\\ mellom to størrelser
som er omvendt proporsjonale.
}
Argumenter for svarene dine.
\fig{1pv24d1opg3}

\grubeksoprecc{1pv24d1opg4}{1PV24D1}
På tørt føre kan bremselengden $B$ (målt i meter) til en bil tilnærmes ved formelen
\[ B=\frac{1}{2}\left(\frac{v}{10}\right)^2\]
hvor $v$ er farten (målt i km/h) bilen hadde før bremsingen startet.
\abc{
\item Vis at en bil som kjører i 70\,km/h har en bremselengde på 24,5 meter.
\item Finn farten når bremselengden er 40,5 meter.
}

\grubeksoprecc{1ph24d1opg2}{1PH24D1}
Forskere har kommet fram til at det er omtrent 20 billiarder maur på jorden.
En billiard er tusen millioner millioner.
\abc{
\item Skriv 20 billiarder på standardform.
\item I en normalt stor maurtue er det mellom 200 000 og 300 000 maur. Anta at en maur
veier mellom 7 mg og 9 mg. Omtrent hvor mange kilogram veier alle maurene i en normalt stor maurtue til sammen?
}

\grubeksoprecc{1ph24d1opg3}{1PH24D1}
Nedenfor er det beskrevet tre situasjoner. Fullfør siste setning i hver beskrivelse ved
å sette inn teksten som står i en av rutene nedenfor. Argumenter for alle tre svarene dine.
\begin{center}
\begin{tabular}{|l|}
	\hline
	proporsjonale
	størrelser \\ \hline
	omvendt
	proporsjonale
	størrelser \\ \hline 
	verken proporsjonale
	eller omvendt
	proporsjonale
	størrelser \\ 
	\hline
\end{tabular}
\end{center} \vsk

\textbf{Situasjon A}\\
Det koster 2200 kroner å leie en badstue. Antallet personer som er med på å betale
leien, og prisen per person er ... \vsk

\textbf{Situasjon B}\\
Når du kjøper brus, kan du ta tre flasker og betale for to. Antallet flasker du kjøper, og
prisen du betaler for alle flaskene, er ... \vsk \vs

\textbf{Situasjon C}\\
Antallet porsjoner vaffelrøre du lager, og mengden mel du trenger, er ...


\eksop{1ph24d2opg2}{1PH24D2}
En lege rører ut et pulver i vann for å lage medisin til en pasient.
Han bruker 6 mg av pulveret per milliliter vann.
Pasienten veier 75 kg, og skal ha 15 mg pulver per kilogram kroppsvekt hvert døgn,
fordelt på tre like store doser. \os

Hvor mange milliliter av medisinen skal pasienten ha i hver dose?

\grubeksoprecc{1ph25d1opg1}{1PV25D1}
Sofie tok buss fra Mandal til Oslo. Bussen holdt en gjennomsnittsfart
på 80\,km/h og brukte 4 timer og 30 minutter på strekningen. \os

Hvor lang er denne strekningen?

\grubeksoprecc{1pv25d1opg5}{1PV25D1}
I blodet er det tre hovedtyper blodceller. De tre hovedtypene er hvite blodceller, røde
blodceller og blodplater. \os

I en liter blod er det $7\cdot10^9$ hvite blodceller, $5\cdot10^{12}$ røde blodceller og $3\cdot10^{11}$ blodplater. \os

Hvor mange blodceller blir dette til sammen? Skriv svaret på standardform.

\grubeksoprecc{1ph25d2opg4}{1PH25D2}
Synnøve er på ferie i USA. En dag går hun innom en butikk for å kjøpe epler.
Hun betaler 4,18 amerikanske dollar for en pose med 3 pund epler.
\begin{itemize}
	\item Pund (lb) er en måleenhet for masse som er vanlig å bruke i USA.
	$1\enh{lb} \approx 0,454\enh{kg} $
	\item Sist Synnøve sjekket valutakursen, tilsvarte 1 amerikansk dollar
	10,16 norske kroner.
\end{itemize}
Hvor mange norske kroner kostet ett kilogram epler?
\end{document}


